%!TEX root = ../dokumentation.tex

\chapter{Theoretical Foundations}
\label{ch:foundations}

\section{Enterprise Knowledge Retrieval}

Enterprise knowledge retrieval encompasses the challenge of surfacing relevant, context-aware information from collections of heterogeneous documents distributed across organisational systems. Modern enterprises accumulate knowledge across structurally distinct source types—issue trackers, collaborative wikis, and document repositories—each of which imposes different constraints on how information is stored, queried, and consumed. The resulting heterogeneity is not merely a matter of file format; it reflects fundamental differences in document length, information density, structural formality, and the semantics of the queries users direct at each source.

The three source types employed in this thesis illustrate the range of structural variation encountered in practice. Issue tracker entries, represented here by GitHub Issues as a proxy for Jira tickets, are short and metadata-rich, with discrete fields for status, labels, and assignees alongside free-text comment threads. Collaborative wiki pages, represented by Markdown documentation as a proxy for Confluence, are discursive and medium-length, organised by human-authored headings into flowing prose. Formal document repositories, represented by PDF whitepapers as a proxy for SharePoint, contain long, hierarchically structured texts with dense technical content and standardised sections. A retrieval approach that treats these source types uniformly risks systematic degradation on at least one category, as enterprise RAG deployments consistently demonstrate when a uniform retrieval configuration produces uneven results across structurally distinct document classes \cite{gao2023ragsurvey}.

Latency expectations compound the difficulty of the retrieval problem. Enterprise users tolerate response times that vary with query complexity: simple factual lookups — retrieving the status of a known issue or locating a specific configuration parameter — carry an implicit expectation of near-instantaneous response, whereas cross-document synthesis queries, which require combining evidence from multiple sources, are accepted with longer wait times. A retrieval system must therefore balance thoroughness against speed in a manner that is sensitive to query type, not only to document volume.

Traditional enterprise search has long relied on lexical matching methods, of which the BM25 ranking function is the canonical representative. BM25 scores documents by term frequency and inverse document frequency, yielding high precision for queries that share surface tokens with the target document. Its limitations emerge when the query and the relevant passage share meaning without sharing vocabulary: a query for ``deployment instructions'' will fail to retrieve a document whose relevant passage reads ``steps to install and configure the service,'' because no term overlap exists. This vocabulary mismatch problem is a persistent source of recall loss in operational retrieval systems.

Federated search addresses a different dimension of the enterprise retrieval problem by coordinating queries across multiple disparate systems and aggregating the results into a unified response. It solves the organisational challenge of querying siloed repositories without requiring data migration, but it does not address the semantic gap: each constituent system still applies lexical matching internally, and the aggregation layer has no mechanism for resolving conceptual relationships across sources. Federated search thus increases coverage without improving the quality of individual matches, leaving cross-source synthesis queries — those that require reconciling information distributed across an issue thread, a wiki page, and a whitepaper — beyond its reach.

The limitations of both approaches converge on a shared root cause: neither represents documents or queries as continuous semantic objects. Keyword search operates on discrete tokens; federated search multiplies the number of token-matching engines without changing the underlying representation. Semantic retrieval methods, which map documents and queries into dense vector representations where proximity reflects meaning rather than surface form, address this root cause directly. The two paradigms examined in this thesis — classical RAG and agent-based retrieval — both build on semantic representations as their foundation, and the comparative evaluation consequently begins from a baseline that lexical search cannot achieve.


\section{Retrieval-Augmented Generation}

Retrieval-Augmented Generation extends the capabilities of large language models (LLMs) by grounding generation in dynamically retrieved external knowledge, rather than relying solely on parameters encoded during training \cite{lewis2020rag}. The original formulation by \citeasnoun{lewis2020rag} decomposes the system into two learnable components: a retriever, parameterised as a Dense Passage Retriever (DPR) that estimates the probability of a document given a query, and a generator, parameterised as a sequence-to-sequence model that conditions its output on both the query and the retrieved passages. A critical architectural property of this design is that the non-parametric retrieval index can be updated or expanded without retraining the generator, making the system's knowledge base directly inspectable and modifiable \cite{lewis2020rag}.

The core RAG pipeline proceeds across three sequential phases: indexing, retrieval, and generation. During indexing, source documents are segmented into chunks, encoded by an embedding model into dense vector representations, and stored in a vector index alongside their source metadata. At query time, the retrieval phase encodes the user query using the same embedding model, computes similarity against the stored vectors — typically using cosine similarity or inner product — and returns the top-$k$ most relevant chunks. These retrieved passages are then concatenated with the original query and passed to the generator as a prompt context, from which the final response is produced. The generator therefore operates on a composed input that merges parametric knowledge with non-parametric evidence, reducing the risk of hallucination on knowledge-intensive tasks.

The RAG landscape has since diversified substantially, and \citeasnoun{gao2023ragsurvey} organise this development into three paradigms: Naive RAG, Advanced RAG, and Modular RAG. Naive RAG implements the core pipeline described above without modification — a single retrieval step, no query transformation, and no post-retrieval filtering — and serves as the tractable baseline against which improvements are measured. However, Naive RAG exhibits well-documented failure modes: retrieval precision is often low, with irrelevant chunks included in the context, while recall gaps mean that relevant passages may not surface at all \cite{gao2023ragsurvey}. These deficiencies propagate directly into generation, where the presence of contradictory or off-topic retrieved content can cause the model to produce responses that conflict with the retrieved evidence.

Advanced RAG addresses these failure modes by introducing pre-retrieval and post-retrieval processing stages around the core pipeline. Pre-retrieval techniques — including query rewriting, query expansion, and Hypothetical Document Embedding (HyDE), where the model generates a synthetic answer used as the retrieval query — improve the semantic alignment between the query and the indexed content. Post-retrieval techniques, most notably reranking with a cross-encoder model, reorder the initial top-$k$ candidates by scoring each passage against the query jointly rather than independently, yielding substantially better relevance ordering \cite{gao2023ragsurvey}. Context compression further reduces noise by truncating or filtering retrieved chunks to the most answer-relevant spans before they are passed to the generator. The classical pipeline implemented in this thesis is positioned as an Advanced RAG system: it performs query-level embedding, vector retrieval, cross-encoder reranking, and context compression, but applies a single fixed retrieval step without source routing or adaptive retrieval decisions.

Modular RAG deconstructs the pipeline further into interchangeable components and introduces new capabilities not present in either earlier paradigm. Routing directs different query types to different retrieval strategies or data sources; scheduling determines whether retrieval is needed at all for a given query; and adaptive retrieval enables the system to retrieve multiple times, conditioning each subsequent query on prior observations \cite{gao2023ragsurvey}. These capabilities bring Modular RAG into conceptual proximity with agent-based systems, in which source routing and adaptive retrieval are externalised as explicit reasoning steps rather than embedded as pipeline parameters.

The effectiveness of a RAG system depends considerably on chunking strategy, which determines the granularity and coherence of the indexed units. Fixed-size chunking is computationally cheap but routinely splits passages at semantically arbitrary boundaries, degrading retrieval precision. Sentence-level chunking respects natural language units yet produces short chunks that may lack sufficient context for accurate relevance scoring. The Small2Big — or hierarchical — strategy offers the strongest balance: small chunks are indexed and retrieved for retrieval precision, but the generator receives the surrounding parent chunk as context, preserving local coherence \cite{gao2023ragsurvey}. Sliding window chunking with configurable overlap reduces the risk of splitting relevant spans across adjacent chunks and remains a practical complement to sentence-level segmentation.

Embedding model choice determines the representational capacity of the retrieval index. Sparse retrieval methods, such as BM25, operate on exact lexical overlap and serve as strong baselines for short, precise queries, but generalise poorly to paraphrase or synonym variation. Dense retrieval models — including the BGE and E5 families, evaluated comparatively on the Massive Text Embedding Benchmark (MTEB) — learn semantic representations that capture meaning beyond surface form, enabling retrieval across queries that share intent but not vocabulary \cite{lewis2020rag}. In practice, hybrid retrieval combining sparse and dense signals can outperform either method alone on heterogeneous corpora.

Despite these enhancements, Advanced RAG remains a static pipeline: retrieval is always performed, always from a single index, and always at a fixed point in the generation process. This design assumption becomes a liability when the information need is ambiguous, when multiple retrieval steps are required to integrate evidence across documents, or when retrieval is unnecessary and introduces noise. \citeasnoun{asai2023selfrag} demonstrate that indiscriminate retrieval, regardless of whether the query requires it, reduces model performance; removing adaptive retrieval decisions from their system causes accuracy to fall from 54.9 to 42.6 on PopQA. \citeasnoun{jiang2023flare} show that a single retrieval step is insufficient for long-form generation tasks where information needs evolve across the output, and that active, forward-looking retrieval outperforms a fixed single-step baseline by a substantial margin on multi-hop reasoning tasks. Taken together, these findings suggest that the architectural ceiling of static retrieval pipelines motivates a fundamentally different approach — one in which retrieval decisions are made dynamically by a system capable of selecting sources, composing multi-step queries, and adjusting its strategy based on intermediate results.


\section{Agent-Based Information Retrieval}

Agent-based systems extend language model capabilities by embedding the model within a decision-making cycle that interleaves reasoning, action, and environmental feedback. Where a classical RAG pipeline executes a fixed sequence — retrieve once, then generate — an agent-based system treats retrieval as one of many possible actions, selected dynamically based on the current query state and prior observations. \citeasnoun{xi2023rise} define this class of system through four distinguishing properties: autonomy, in that the agent operates without direct human intervention at each step; reactivity, in that it responds to observations returned by tools; pro-activeness, in that it pursues a goal across multiple steps; and social ability, in that it can in principle coordinate with other agents or human users. These properties collectively describe a system that decides \emph{how} to retrieve information, rather than one that simply retrieves it.

The ReAct framework, introduced by \citeasnoun{yao2023react}, provides the theoretical backbone for this class of agent by interleaving verbal reasoning traces with discrete tool actions in a repeating loop. Each cycle consists of three elements: a Thought, which is a free-form reasoning step in which the model assesses its current state and decides on a next action; an Action, which is a structured call to a defined tool with specific arguments; and an Observation, which is the value returned by that tool and appended to the context before the next Thought step. The agent continues cycling until it determines that sufficient information has been gathered to produce a final answer. \citeasnoun{yao2023react} demonstrate that this interleaving significantly outperforms either reasoning or acting in isolation, reporting 27.4 exact-match accuracy on HotpotQA and 71\% task success on ALFWorld. Critically, the Thought traces are human-readable, making the agent's decision process auditable in a way that the single-pass RAG pipeline is not.

Tool use is the mechanism through which the agent grounds its reasoning in external information sources, and it directly addresses several limitations of plain language model inference. \citeasnoun{schick2023toolformer} identify four such limitations: the model's knowledge is bounded by its training cut-off; in-context missteering can cause factually incorrect completions; hallucination degrades reliability on knowledge-intensive tasks; and domain-specific sub-tasks such as arithmetic require capabilities the base model lacks. The filtering criterion that distinguishes a useful tool call from a redundant one is whether the tool's response reduces future generation loss by at least a fixed threshold — a principle that generalises beyond Toolformer's specific implementation to any agent that must decide whether to query an external source \cite{schick2023toolformer}. The ReAct loop overcomes the sequential dependency limitation inherent in single-step tool use by making each Observation available as context for the subsequent Thought, enabling the agent to compose tool calls iteratively across a multi-hop query.

\citeasnoun{wang2023survey} decompose agent architecture into four interacting modules that map directly onto the system built in this thesis. The profiling module defines the agent's role and constraints, corresponding to the system prompt that instructs the agent to function as an enterprise retrieval assistant with access to source-specific tools. The memory module encompasses both short-term working memory — the in-context accumulation of the running ReAct trace — and long-term semantic memory, implemented here as the per-source document indexes. The planning module governs how the agent sequences its actions; \citeasnoun{wang2023survey} distinguish single-path planning (Chain-of-Thought), multi-path planning (Tree-of-Thought), and feedback-driven planning (ReAct), with ReAct identified as optimal for multi-hop retrieval tasks where each step must inform the next. The action module defines the typed tool calls available to the agent, corresponding in this thesis to three source-specific retrieval tools: one for GitHub Issues, one for Markdown documentation, and one for PDF whitepapers.

The Cognitive Architectures for Language Agents (CoALA) framework from \citeasnoun{sumers2023coala} provides a formal basis for distinguishing agent-based systems from classical RAG. CoALA partitions agent memory into working memory (the current query, retrieved passages, and running trace), episodic memory (past trajectories), semantic memory (long-term unstructured knowledge), and procedural memory (implicit weights plus explicit prompt templates and tool definitions). Under this taxonomy, classical RAG is a constrained agent that reads from a fixed, read-only semantic memory with no surrounding decision cycle — \citeasnoun{sumers2023coala} describe retrieval-augmented methods as ``retrieving from a semantic memory of unstructured text'' and position classical RAG as the degenerate case in which the agent has no planning stage and no adaptive retrieval decision. The thesis agent adds a decision cycle around retrieval, enabling conditional and iterative access to multiple semantic memory stores rather than a single unified index.

\citeasnoun{shinn2023reflexion} extend the ReAct loop with verbal self-reflection, in which a dedicated self-reflection model generates a textual summary of why a prior attempt failed, storing that summary in episodic memory for the next attempt. The resulting system achieves 97\% task success on ALFWorld, compared to 78\% for standard ReAct, demonstrating that verbal reinforcement of the decision cycle substantially improves agent performance on complex sequential tasks. This capability is not implemented in the agent evaluated in this thesis, as queries are treated as independent and no learning is persisted between runs; Reflexion nonetheless represents a natural extension direction that would bring the system closer to a continuously adaptive enterprise retrieval agent.

The contrast between classical RAG and agent-based retrieval is most sharply visible across the evaluation dimensions examined in this thesis. Classical RAG routes all queries through a single unified index, retrieves once per query, and produces a response in a single language model call, resulting in lower latency and lower token expenditure. The agent-based system routes each query to one or more source-specific tools, retrieves conditionally and iteratively based on accumulated observations, and may invoke the language model several times before producing a final response. This iterative structure improves support for multi-hop queries that require synthesising information from multiple documents or sources, but imposes measurable costs in latency and token consumption — a trade-off \citeasnoun{gao2023ragsurvey} identify as central to the design space of advanced retrieval systems. The readable Thought traces produced by the ReAct loop partially offset the higher cost by providing interpretability that the black-box RAG pipeline cannot match.


\section{Evaluation Methodology}

Evaluating retrieval-augmented systems presents a challenge that traditional natural language processing (NLP) benchmarks are not designed to address: the quality of a generated answer depends jointly on retrieval precision and generation faithfulness, and these two failure modes require distinct measurement instruments. The evaluation framework employed in this thesis combines three complementary layers—reference-free retrieval-grounded metrics, reference-dependent semantic similarity measures, and non-quality operational metrics—to produce a comprehensive and fair comparison between the classical RAG pipeline and the agent-based system.

\subsection*{RAGAS: Reference-Free Retrieval-Grounded Evaluation}

The RAGAS framework \cite{es2023ragas} provides three metrics that assess the pipeline end-to-end without requiring pre-existing ground truth answer annotations. Context Relevance measures whether retrieved passages are genuinely necessary for answering the query. An LLM is prompted to extract only those sentences from the retrieved context that are essential for a response, yielding a ratio $\text{CR} = |S_{\text{extracted}}| / |S_{\text{total}}|$, where a score approaching 1.0 indicates precise, low-noise retrieval. This metric is the most sensitive to retrieval strategy differences between the two systems, though it also shows the weakest alignment with human judgements (0.70 pairwise accuracy), partly due to extraction failures on long passages \cite{es2023ragas}.

Answer Faithfulness measures the degree to which a generated answer is grounded in the retrieved context, rather than in parametric model knowledge. The metric decomposes the answer into atomic claims, verifies each binary claim against the context, and computes $\text{Faithfulness} = |V| / |S|$, where $V$ denotes verified claims and $S$ the total claim count. A score of 1.0 indicates full contextual grounding, while lower scores signal hallucination relative to the retrieved evidence. Validation on the WikiEval benchmark demonstrates 0.95 pairwise accuracy against human judgements, making Faithfulness the most reliable RAGAS metric \cite{es2023ragas}.

Answer Relevance captures whether the generated response addresses the original query, penalising verbose, incomplete, or off-topic answers. An LLM generates $N$ synthetic paraphrases of the query from the answer text, and the metric computes $\text{AR} = \frac{1}{N} \sum_{i=1}^{N} \cos(\mathbf{e}(q), \mathbf{e}(\hat{q}_i))$, where $\mathbf{e}(\cdot)$ denotes the embedding function. Because RAGAS operates without reference answers, it supports evaluation on synthetic corpora where no pre-existing annotations exist \cite{es2023ragas}. However, RAGAS cannot detect source routing errors: an answer that is fully grounded in a retrieved passage from the wrong source type still achieves a Faithfulness score of 1.0. Source routing accuracy must therefore be measured as a separate, agent-specific metric.

\subsection*{BERTScore: Reference-Dependent Semantic Similarity}

Lexical overlap metrics such as BLEU and ROUGE-L were designed for machine translation and summarisation tasks in which the reference text and the candidate response share significant surface form. In a retrieval setting where correct answers may be expressed through diverse valid paraphrases, BLEU produces near-random correlation with human judgements—reaching a Pearson correlation of $-0.019$ on the COCO image captioning benchmark \cite{zhang2020bertscore}—and therefore provides insufficient sensitivity for this evaluation.

BERTScore \cite{zhang2020bertscore} addresses this gap by computing soft token-level matching between candidate and reference using contextual embeddings from RoBERTa-Large. Recall, precision, and F1 are defined as:
\begin{align}
R_{\text{BERT}} &= \frac{1}{|x|} \sum_{i} \max_{j}\, \mathbf{x}_i^\top \hat{\mathbf{x}}_j \\
P_{\text{BERT}} &= \frac{1}{|\hat{x}|} \sum_{j} \max_{i}\, \mathbf{x}_i^\top \hat{\mathbf{x}}_j \\
F_{\text{BERT}} &= 2 \cdot \frac{P_{\text{BERT}} \cdot R_{\text{BERT}}}{P_{\text{BERT}} + R_{\text{BERT}}}
\end{align}
where $\mathbf{x}_i$ and $\hat{\mathbf{x}}_j$ are the contextual token embeddings of the reference and candidate respectively, optionally weighted by inverse document frequency to down-weight function words. BERTScore achieves a recall correlation of 0.917 under IDF weighting on the same captioning benchmark, and its advantage over BLEU is largest when paraphrase diversity is high \cite{zhang2020bertscore}—exactly the condition that characterises open-domain RAG evaluation.

BERTScore and RAGAS are complementary rather than redundant. RAGAS Faithfulness is reference-free and measures grounding in retrieved context; BERTScore is reference-dependent and measures semantic proximity to a correct ground truth answer. An answer grounded faithfully in incomplete retrieved content may score highly on Faithfulness yet diverge substantially from the correct answer, a gap that only BERTScore captures. Together, the two instruments cover both retrieval-side and generation-side quality without redundancy.

\subsection*{Four Capability Dimensions from Gao et al.}

Beyond aggregate metric scores, \citeasnoun{gao2023ragsurvey} identify four robustness dimensions that characterise the qualitative behaviour of retrieval-augmented systems under realistic operating conditions. Noise robustness concerns whether a system can identify and use the relevant passages within a context that also contains irrelevant or contradictory material. Negative rejection tests whether a system correctly declines to answer when no retrieved passage contains useful information, rather than hallucinating a plausible-sounding response. Information integration measures whether the system can synthesise content from multiple retrieved passages to answer multi-hop queries, a capability that single-pass RAG struggles with by design. Counterfactual robustness assesses whether a system can identify and discount factual errors present in retrieved passages. These four dimensions directly inform the query taxonomy used in this evaluation: simple factual queries probe baseline retrieval quality, cross-document synthesis queries test information integration, and ambiguous or incomplete queries probe negative rejection and noise robustness.

\subsection*{Synthetic Query Generation and LLM-as-Judge}

No pre-existing annotated query set covers the specific proxy corpus—GitHub Issues, Markdown documentation, and PDF whitepapers—used in this thesis. Synthetic QA generation addresses this constraint within a zero-budget academic context: an LLM is prompted to produce candidate question–answer pairs from sampled document passages, and the resulting set is manually reviewed and filtered against quality criteria before use. RAGAS explicitly validates this methodology, demonstrating that reference-free metrics achieve 0.70–0.95 pairwise accuracy with human judgements and that the framework ``enables faster evaluation cycles of RAG architectures without relying on ground truth annotations'' \cite{es2023ragas}. The manual review step is essential for academic credibility, as it filters incoherent, unanswerable, or trivially simple items that automated generation produces. Ground truth answers produced during this process additionally serve as the reference for BERTScore computation.

\subsection*{Operational and Agent-Specific Metrics}

Quality metrics alone are insufficient for an architectural comparison that must inform practical deployment decisions. Latency is measured as wall-clock time from query submission to complete response, averaged over multiple runs, and includes all LLM API calls and, for the agent, all intermediate tool execution steps. Token cost is derived from input and output token counts in API response objects, converted to an estimated cost per query and per 1,000 queries; the agent incurs structurally higher token consumption due to multi-step reasoning traces. Two agent-specific metrics are tracked in addition: tool call count per query (average and stratified by query type) and source routing accuracy, defined as the proportion of queries for which the agent queries the correct source type relative to a manually labelled ground truth. Implementation complexity is assessed through a structured rubric covering flexibility, debuggability, and extensibility, proxied by objective measures including lines of code per component, number of configuration parameters, and the time required to onboard an additional data source. \citeasnoun{ferrazzi2026agenticrag} explicitly exclude implementation complexity from their evaluation scope; its inclusion here as a formal dimension constitutes a direct contribution of this thesis.

\begin{table}[ht]
\centering
\caption{Evaluation metrics employed in this thesis, organised by dimension.}
\label{tab:metrics}
\begin{tabular}{lll}
\hline
\textbf{Dimension} & \textbf{Metric} & \textbf{Reference-free?} \\
\hline
Context relevance & RAGAS Context Relevance & Yes \\
Answer faithfulness & RAGAS Faithfulness & Yes \\
Answer relevance & RAGAS Answer Relevance & Yes \\
Semantic similarity & BERTScore F$_1$ (RoBERTa-Large) & No \\
Latency & Mean wall-clock time per query & --- \\
Token cost & API token counts $\rightarrow$ estimated USD & --- \\
Agent efficiency & Tool calls per query & --- \\
Agent routing & Source selection accuracy & No \\
Complexity & Rubric score + objective proxies & --- \\
\hline
\end{tabular}
\end{table}

Taken together, these considerations suggest that no single metric is sufficient to characterise the difference between a classical RAG pipeline and an agent-based system, and that a multi-dimensional framework spanning retrieval quality, generation fidelity, operational cost, and structural complexity is necessary to produce findings that are both academically rigorous and practically actionable.
